% Include guard: safe to % Include guard: safe to % Include guard: safe to % Include guard: safe to \input{macros} more than once (e.g. once directly
% and once via common.tex, which also loads it).
\ifdefined\pfplMacrosLoaded\endinput\fi
\global\let\pfplMacrosLoaded\empty

% --- Sequence notation ---
% Expand an expression with one free variable s' (e.g. s(s')) into the
% sequence obtained by ranging s' over 1..n: \seqapp{s} => s(1), \dots, s(n)
% Optional second argument overrides the delimiter (default ","):
%   \seqapp{s}[;] => s(1); \dots; s(n)
\RequirePackage{xparse}
\RequirePackage{textgreek}
\NewDocumentCommand{\seqapp}{ m O{,} }
 {
   #1(1)#2\dots#2#1(n)
 }

% Like \seqapp, but #1 is the NAME (no backslash) of an existing one-argument
% macro to invoke at index 1 and index n, e.g. \seqcall{indexedValence} =>
% \indexedValence{1}, \dots, \indexedValence{n}. Optional delimiter as above.
\NewDocumentCommand{\seqcall}{ m O{,} }
 {
   \csname #1\endcsname{1}#2\dots#2\csname #1\endcsname{n}
 }

% --- Math-alphabet shorthand generator ---
% \GenMathShorthand{prefix}{\mathXX}{letters} defines, for each letter L in
% `letters`, \prefixL := \mathXX{L}. `letters` is just typed out (e.g.
% ABC...Z), not a range expression, so any subset/order is fine.
\RequirePackage{xparse}
\ExplSyntaxOn
\NewDocumentCommand{\GenMathShorthand}{ m m m }
 {
   \tl_map_inline:nn { #3 }
     {
       \exp_args:Ncx \newcommand { #1 ##1 } { \exp_not:N #2 {##1} }
     }
 }
\ExplSyntaxOff

% --- Format-composition dispatcher ---
% \NewFmt{c}{\cmd} registers a format character c (any single token/char) so
% it can appear in \f's format-code prefix as a stand-in for \cmd, which must
% take exactly one argument.
% \f <fmt>.<word> composes the registered commands for each character of
% <fmt>, left to right = outermost to innermost, around <word>:
%   \f vc.A     => \vec{\mathcal{A}}
%   \f b.{xy}   => \mathbf{xy}          (braced word for more than one token)
% Text-mode commands (e.g. \texttt, registered as 't') must come last/
% innermost when mixed with math-only commands (\vec, \mathcal, \mathbf,
% \mathit, ...): \f it.A works (\mathit{\texttt{A}}), \f ti.A does not
% (\texttt switches to text/LR mode, and \mathit then fails outside math).
\ExplSyntaxOn
\prop_new:N \g__pfpl_fmt_prop
\tl_new:N \l__pfpl_f_fmt_tl
\tl_new:N \l__pfpl_f_acc_tl

\msg_new:nnn { pfpl } { fmt-undefined }
  { Format~character~'#1'~is~not~defined;~use~\iow_char:N\\NewFmt~to~register~it~first. }
\msg_new:nnn { pfpl } { fmt-duplicate }
  { Format~character~'#1'~is~already~defined. }

\NewDocumentCommand{\NewFmt}{ m m }
 {
   \prop_if_in:NnTF \g__pfpl_fmt_prop {#1}
     { \msg_error:nnn { pfpl } { fmt-duplicate } {#1} }
     { \prop_gput:Nnn \g__pfpl_fmt_prop {#1} {#2} }
 }

\NewDocumentCommand{\f}{ u{.} m }
 {
   \group_begin:
     \tl_set:Nn \l__pfpl_f_fmt_tl {#1}
     \tl_reverse:N \l__pfpl_f_fmt_tl
     \tl_set:Nn \l__pfpl_f_acc_tl {#2}
     \tl_map_inline:Nn \l__pfpl_f_fmt_tl
       {
         \prop_if_in:NnTF \g__pfpl_fmt_prop {##1}
           {
             \tl_set:Nx \l__pfpl_f_acc_tl
               {
                 \prop_item:Nn \g__pfpl_fmt_prop {##1}
                   { \exp_not:V \l__pfpl_f_acc_tl }
               }
           }
           { \msg_error:nnn { pfpl } { fmt-undefined } {##1} }
       }
     \tl_use:N \l__pfpl_f_acc_tl
   \group_end:
 }
\ExplSyntaxOff

\NewFmt{v}{\vec}
\NewFmt{c}{\mathcal}
\NewFmt{b}{\mathbf}
\NewFmt{i}{\mathit}
\NewFmt{t}{\texttt}

% --- Single-letter shorthands (via \GenMathShorthand) ---
% \vA, \vB, ..., \vz := \vec{A}, \vec{B}, ..., \vec{z}
% \cA, \cB, ..., \cZ := \mathcal{A}, \mathcal{B}, ..., \mathcal{Z}
% Uses plain \newcommand internally, so any name collision errors loudly
% instead of silently overriding an existing command.
\GenMathShorthand{v}{\vec}{ABCDEFGHIJKLMNOPQRSTUVWXYZabcdefghijklmnopqrstuvwxyz}
\GenMathShorthand{c}{\mathcal}{ABCDEFGHIJKLMNOPQRSTUVWXYZ}

% --- Alias generator ---
% \GenAlias{newprefix}{oldprefix}{names} defines, for each comma-separated
% name in `names`, \newprefixNAME as an alias (\let) for \oldprefixNAME.
% Unlike \GenMathShorthand, `names` can be multi-character (e.g. Greek letter
% names), and the target must already exist and take no arguments (\let just
% copies its meaning, it doesn't wrap it). Errors loudly instead of silently
% overriding if \newprefixNAME is already defined.
\ExplSyntaxOn
\msg_new:nnn { pfpl } { alias-duplicate }
  { Shorthand~\iow_char:N\\#1~already~defined;~not~aliased. }

\NewDocumentCommand{\GenAlias}{ m m m }
 {
   \clist_map_inline:nn {#3}
     {
       \exp_args:Nc \cs_if_exist:NTF { #1 ##1 }
         { \msg_error:nnn { pfpl } { alias-duplicate } { #1 ##1 } }
         { \exp_args:Ncc \let { #1 ##1 } { #2 ##1 } }
     }
 }
\ExplSyntaxOff
 more than once (e.g. once directly
% and once via common.tex, which also loads it).
\ifdefined\pfplMacrosLoaded\endinput\fi
\global\let\pfplMacrosLoaded\empty

% --- Sequence notation ---
% Expand an expression with one free variable s' (e.g. s(s')) into the
% sequence obtained by ranging s' over 1..n: \seqapp{s} => s(1), \dots, s(n)
% Optional second argument overrides the delimiter (default ","):
%   \seqapp{s}[;] => s(1); \dots; s(n)
\RequirePackage{xparse}
\RequirePackage{textgreek}
\NewDocumentCommand{\seqapp}{ m O{,} }
 {
   #1(1)#2\dots#2#1(n)
 }

% Like \seqapp, but #1 is the NAME (no backslash) of an existing one-argument
% macro to invoke at index 1 and index n, e.g. \seqcall{indexedValence} =>
% \indexedValence{1}, \dots, \indexedValence{n}. Optional delimiter as above.
\NewDocumentCommand{\seqcall}{ m O{,} }
 {
   \csname #1\endcsname{1}#2\dots#2\csname #1\endcsname{n}
 }

% --- Math-alphabet shorthand generator ---
% \GenMathShorthand{prefix}{\mathXX}{letters} defines, for each letter L in
% `letters`, \prefixL := \mathXX{L}. `letters` is just typed out (e.g.
% ABC...Z), not a range expression, so any subset/order is fine.
\RequirePackage{xparse}
\ExplSyntaxOn
\NewDocumentCommand{\GenMathShorthand}{ m m m }
 {
   \tl_map_inline:nn { #3 }
     {
       \exp_args:Ncx \newcommand { #1 ##1 } { \exp_not:N #2 {##1} }
     }
 }
\ExplSyntaxOff

% --- Format-composition dispatcher ---
% \NewFmt{c}{\cmd} registers a format character c (any single token/char) so
% it can appear in \f's format-code prefix as a stand-in for \cmd, which must
% take exactly one argument.
% \f <fmt>.<word> composes the registered commands for each character of
% <fmt>, left to right = outermost to innermost, around <word>:
%   \f vc.A     => \vec{\mathcal{A}}
%   \f b.{xy}   => \mathbf{xy}          (braced word for more than one token)
% Text-mode commands (e.g. \texttt, registered as 't') must come last/
% innermost when mixed with math-only commands (\vec, \mathcal, \mathbf,
% \mathit, ...): \f it.A works (\mathit{\texttt{A}}), \f ti.A does not
% (\texttt switches to text/LR mode, and \mathit then fails outside math).
\ExplSyntaxOn
\prop_new:N \g__pfpl_fmt_prop
\tl_new:N \l__pfpl_f_fmt_tl
\tl_new:N \l__pfpl_f_acc_tl

\msg_new:nnn { pfpl } { fmt-undefined }
  { Format~character~'#1'~is~not~defined;~use~\iow_char:N\\NewFmt~to~register~it~first. }
\msg_new:nnn { pfpl } { fmt-duplicate }
  { Format~character~'#1'~is~already~defined. }

\NewDocumentCommand{\NewFmt}{ m m }
 {
   \prop_if_in:NnTF \g__pfpl_fmt_prop {#1}
     { \msg_error:nnn { pfpl } { fmt-duplicate } {#1} }
     { \prop_gput:Nnn \g__pfpl_fmt_prop {#1} {#2} }
 }

\NewDocumentCommand{\f}{ u{.} m }
 {
   \group_begin:
     \tl_set:Nn \l__pfpl_f_fmt_tl {#1}
     \tl_reverse:N \l__pfpl_f_fmt_tl
     \tl_set:Nn \l__pfpl_f_acc_tl {#2}
     \tl_map_inline:Nn \l__pfpl_f_fmt_tl
       {
         \prop_if_in:NnTF \g__pfpl_fmt_prop {##1}
           {
             \tl_set:Nx \l__pfpl_f_acc_tl
               {
                 \prop_item:Nn \g__pfpl_fmt_prop {##1}
                   { \exp_not:V \l__pfpl_f_acc_tl }
               }
           }
           { \msg_error:nnn { pfpl } { fmt-undefined } {##1} }
       }
     \tl_use:N \l__pfpl_f_acc_tl
   \group_end:
 }
\ExplSyntaxOff

\NewFmt{v}{\vec}
\NewFmt{c}{\mathcal}
\NewFmt{b}{\mathbf}
\NewFmt{i}{\mathit}
\NewFmt{t}{\texttt}

% --- Single-letter shorthands (via \GenMathShorthand) ---
% \vA, \vB, ..., \vz := \vec{A}, \vec{B}, ..., \vec{z}
% \cA, \cB, ..., \cZ := \mathcal{A}, \mathcal{B}, ..., \mathcal{Z}
% Uses plain \newcommand internally, so any name collision errors loudly
% instead of silently overriding an existing command.
\GenMathShorthand{v}{\vec}{ABCDEFGHIJKLMNOPQRSTUVWXYZabcdefghijklmnopqrstuvwxyz}
\GenMathShorthand{c}{\mathcal}{ABCDEFGHIJKLMNOPQRSTUVWXYZ}

% --- Alias generator ---
% \GenAlias{newprefix}{oldprefix}{names} defines, for each comma-separated
% name in `names`, \newprefixNAME as an alias (\let) for \oldprefixNAME.
% Unlike \GenMathShorthand, `names` can be multi-character (e.g. Greek letter
% names), and the target must already exist and take no arguments (\let just
% copies its meaning, it doesn't wrap it). Errors loudly instead of silently
% overriding if \newprefixNAME is already defined.
\ExplSyntaxOn
\msg_new:nnn { pfpl } { alias-duplicate }
  { Shorthand~\iow_char:N\\#1~already~defined;~not~aliased. }

\NewDocumentCommand{\GenAlias}{ m m m }
 {
   \clist_map_inline:nn {#3}
     {
       \exp_args:Nc \cs_if_exist:NTF { #1 ##1 }
         { \msg_error:nnn { pfpl } { alias-duplicate } { #1 ##1 } }
         { \exp_args:Ncc \let { #1 ##1 } { #2 ##1 } }
     }
 }
\ExplSyntaxOff
 more than once (e.g. once directly
% and once via common.tex, which also loads it).
\ifdefined\pfplMacrosLoaded\endinput\fi
\global\let\pfplMacrosLoaded\empty

% --- Sequence notation ---
% Expand an expression with one free variable s' (e.g. s(s')) into the
% sequence obtained by ranging s' over 1..n: \seqapp{s} => s(1), \dots, s(n)
% Optional second argument overrides the delimiter (default ","):
%   \seqapp{s}[;] => s(1); \dots; s(n)
\RequirePackage{xparse}
\RequirePackage{textgreek}
\NewDocumentCommand{\seqapp}{ m O{,} }
 {
   #1(1)#2\dots#2#1(n)
 }

% Like \seqapp, but #1 is the NAME (no backslash) of an existing one-argument
% macro to invoke at index 1 and index n, e.g. \seqcall{indexedValence} =>
% \indexedValence{1}, \dots, \indexedValence{n}. Optional delimiter as above.
\NewDocumentCommand{\seqcall}{ m O{,} }
 {
   \csname #1\endcsname{1}#2\dots#2\csname #1\endcsname{n}
 }

% --- Math-alphabet shorthand generator ---
% \GenMathShorthand{prefix}{\mathXX}{letters} defines, for each letter L in
% `letters`, \prefixL := \mathXX{L}. `letters` is just typed out (e.g.
% ABC...Z), not a range expression, so any subset/order is fine.
\RequirePackage{xparse}
\ExplSyntaxOn
\NewDocumentCommand{\GenMathShorthand}{ m m m }
 {
   \tl_map_inline:nn { #3 }
     {
       \exp_args:Ncx \newcommand { #1 ##1 } { \exp_not:N #2 {##1} }
     }
 }
\ExplSyntaxOff

% --- Format-composition dispatcher ---
% \NewFmt{c}{\cmd} registers a format character c (any single token/char) so
% it can appear in \f's format-code prefix as a stand-in for \cmd, which must
% take exactly one argument.
% \f <fmt>.<word> composes the registered commands for each character of
% <fmt>, left to right = outermost to innermost, around <word>:
%   \f vc.A     => \vec{\mathcal{A}}
%   \f b.{xy}   => \mathbf{xy}          (braced word for more than one token)
% Text-mode commands (e.g. \texttt, registered as 't') must come last/
% innermost when mixed with math-only commands (\vec, \mathcal, \mathbf,
% \mathit, ...): \f it.A works (\mathit{\texttt{A}}), \f ti.A does not
% (\texttt switches to text/LR mode, and \mathit then fails outside math).
\ExplSyntaxOn
\prop_new:N \g__pfpl_fmt_prop
\tl_new:N \l__pfpl_f_fmt_tl
\tl_new:N \l__pfpl_f_acc_tl

\msg_new:nnn { pfpl } { fmt-undefined }
  { Format~character~'#1'~is~not~defined;~use~\iow_char:N\\NewFmt~to~register~it~first. }
\msg_new:nnn { pfpl } { fmt-duplicate }
  { Format~character~'#1'~is~already~defined. }

\NewDocumentCommand{\NewFmt}{ m m }
 {
   \prop_if_in:NnTF \g__pfpl_fmt_prop {#1}
     { \msg_error:nnn { pfpl } { fmt-duplicate } {#1} }
     { \prop_gput:Nnn \g__pfpl_fmt_prop {#1} {#2} }
 }

\NewDocumentCommand{\f}{ u{.} m }
 {
   \group_begin:
     \tl_set:Nn \l__pfpl_f_fmt_tl {#1}
     \tl_reverse:N \l__pfpl_f_fmt_tl
     \tl_set:Nn \l__pfpl_f_acc_tl {#2}
     \tl_map_inline:Nn \l__pfpl_f_fmt_tl
       {
         \prop_if_in:NnTF \g__pfpl_fmt_prop {##1}
           {
             \tl_set:Nx \l__pfpl_f_acc_tl
               {
                 \prop_item:Nn \g__pfpl_fmt_prop {##1}
                   { \exp_not:V \l__pfpl_f_acc_tl }
               }
           }
           { \msg_error:nnn { pfpl } { fmt-undefined } {##1} }
       }
     \tl_use:N \l__pfpl_f_acc_tl
   \group_end:
 }
\ExplSyntaxOff

\NewFmt{v}{\vec}
\NewFmt{c}{\mathcal}
\NewFmt{b}{\mathbf}
\NewFmt{i}{\mathit}
\NewFmt{t}{\texttt}

% --- Single-letter shorthands (via \GenMathShorthand) ---
% \vA, \vB, ..., \vz := \vec{A}, \vec{B}, ..., \vec{z}
% \cA, \cB, ..., \cZ := \mathcal{A}, \mathcal{B}, ..., \mathcal{Z}
% Uses plain \newcommand internally, so any name collision errors loudly
% instead of silently overriding an existing command.
\GenMathShorthand{v}{\vec}{ABCDEFGHIJKLMNOPQRSTUVWXYZabcdefghijklmnopqrstuvwxyz}
\GenMathShorthand{c}{\mathcal}{ABCDEFGHIJKLMNOPQRSTUVWXYZ}

% --- Alias generator ---
% \GenAlias{newprefix}{oldprefix}{names} defines, for each comma-separated
% name in `names`, \newprefixNAME as an alias (\let) for \oldprefixNAME.
% Unlike \GenMathShorthand, `names` can be multi-character (e.g. Greek letter
% names), and the target must already exist and take no arguments (\let just
% copies its meaning, it doesn't wrap it). Errors loudly instead of silently
% overriding if \newprefixNAME is already defined.
\ExplSyntaxOn
\msg_new:nnn { pfpl } { alias-duplicate }
  { Shorthand~\iow_char:N\\#1~already~defined;~not~aliased. }

\NewDocumentCommand{\GenAlias}{ m m m }
 {
   \clist_map_inline:nn {#3}
     {
       \exp_args:Nc \cs_if_exist:NTF { #1 ##1 }
         { \msg_error:nnn { pfpl } { alias-duplicate } { #1 ##1 } }
         { \exp_args:Ncc \let { #1 ##1 } { #2 ##1 } }
     }
 }
\ExplSyntaxOff
 more than once (e.g. once directly
% and once via common.tex, which also loads it).
\ifdefined\pfplMacrosLoaded\endinput\fi
\global\let\pfplMacrosLoaded\empty

% --- Sequence notation ---
% Expand an expression with one free variable s' (e.g. s(s')) into the
% sequence obtained by ranging s' over 1..n: \seqapp{s} => s(1), \dots, s(n)
% Optional second argument overrides the delimiter (default ","):
%   \seqapp{s}[;] => s(1); \dots; s(n)
\RequirePackage{xparse}
\RequirePackage{textgreek}
\NewDocumentCommand{\seqapp}{ m O{,} }
 {
   #1(1)#2\dots#2#1(n)
 }

% Like \seqapp, but #1 is the NAME (no backslash) of an existing one-argument
% macro to invoke at index 1 and index n, e.g. \seqcall{indexedValence} =>
% \indexedValence{1}, \dots, \indexedValence{n}. Optional delimiter as above.
\NewDocumentCommand{\seqcall}{ m O{,} }
 {
   \csname #1\endcsname{1}#2\dots#2\csname #1\endcsname{n}
 }

% --- Math-alphabet shorthand generator ---
% \GenMathShorthand{prefix}{\mathXX}{letters} defines, for each letter L in
% `letters`, \prefixL := \mathXX{L}. `letters` is just typed out (e.g.
% ABC...Z), not a range expression, so any subset/order is fine.
\RequirePackage{xparse}
\ExplSyntaxOn
\NewDocumentCommand{\GenMathShorthand}{ m m m }
 {
   \tl_map_inline:nn { #3 }
     {
       \exp_args:Ncx \newcommand { #1 ##1 } { \exp_not:N #2 {##1} }
     }
 }
\ExplSyntaxOff

% --- Format-composition dispatcher ---
% \NewFmt{c}{\cmd} registers a format character c (any single token/char) so
% it can appear in \f's format-code prefix as a stand-in for \cmd, which must
% take exactly one argument.
% \f <fmt>.<word> composes the registered commands for each character of
% <fmt>, left to right = outermost to innermost, around <word>:
%   \f vc.A     => \vec{\mathcal{A}}
%   \f b.{xy}   => \mathbf{xy}          (braced word for more than one token)
% Text-mode commands (e.g. \texttt, registered as 't') must come last/
% innermost when mixed with math-only commands (\vec, \mathcal, \mathbf,
% \mathit, ...): \f it.A works (\mathit{\texttt{A}}), \f ti.A does not
% (\texttt switches to text/LR mode, and \mathit then fails outside math).
\ExplSyntaxOn
\prop_new:N \g__pfpl_fmt_prop
\tl_new:N \l__pfpl_f_fmt_tl
\tl_new:N \l__pfpl_f_acc_tl

\msg_new:nnn { pfpl } { fmt-undefined }
  { Format~character~'#1'~is~not~defined;~use~\iow_char:N\\NewFmt~to~register~it~first. }
\msg_new:nnn { pfpl } { fmt-duplicate }
  { Format~character~'#1'~is~already~defined. }

\NewDocumentCommand{\NewFmt}{ m m }
 {
   \prop_if_in:NnTF \g__pfpl_fmt_prop {#1}
     { \msg_error:nnn { pfpl } { fmt-duplicate } {#1} }
     { \prop_gput:Nnn \g__pfpl_fmt_prop {#1} {#2} }
 }

\NewDocumentCommand{\f}{ u{.} m }
 {
   \group_begin:
     \tl_set:Nn \l__pfpl_f_fmt_tl {#1}
     \tl_reverse:N \l__pfpl_f_fmt_tl
     \tl_set:Nn \l__pfpl_f_acc_tl {#2}
     \tl_map_inline:Nn \l__pfpl_f_fmt_tl
       {
         \prop_if_in:NnTF \g__pfpl_fmt_prop {##1}
           {
             \tl_set:Nx \l__pfpl_f_acc_tl
               {
                 \prop_item:Nn \g__pfpl_fmt_prop {##1}
                   { \exp_not:V \l__pfpl_f_acc_tl }
               }
           }
           { \msg_error:nnn { pfpl } { fmt-undefined } {##1} }
       }
     \tl_use:N \l__pfpl_f_acc_tl
   \group_end:
 }
\ExplSyntaxOff

\NewFmt{v}{\vec}
\NewFmt{c}{\mathcal}
\NewFmt{b}{\mathbf}
\NewFmt{i}{\mathit}
\NewFmt{t}{\texttt}

% --- Single-letter shorthands (via \GenMathShorthand) ---
% \vA, \vB, ..., \vz := \vec{A}, \vec{B}, ..., \vec{z}
% \cA, \cB, ..., \cZ := \mathcal{A}, \mathcal{B}, ..., \mathcal{Z}
% Uses plain \newcommand internally, so any name collision errors loudly
% instead of silently overriding an existing command.
\GenMathShorthand{v}{\vec}{ABCDEFGHIJKLMNOPQRSTUVWXYZabcdefghijklmnopqrstuvwxyz}
\GenMathShorthand{c}{\mathcal}{ABCDEFGHIJKLMNOPQRSTUVWXYZ}

% --- Alias generator ---
% \GenAlias{newprefix}{oldprefix}{names} defines, for each comma-separated
% name in `names`, \newprefixNAME as an alias (\let) for \oldprefixNAME.
% Unlike \GenMathShorthand, `names` can be multi-character (e.g. Greek letter
% names), and the target must already exist and take no arguments (\let just
% copies its meaning, it doesn't wrap it). Errors loudly instead of silently
% overriding if \newprefixNAME is already defined.
\ExplSyntaxOn
\msg_new:nnn { pfpl } { alias-duplicate }
  { Shorthand~\iow_char:N\\#1~already~defined;~not~aliased. }

\NewDocumentCommand{\GenAlias}{ m m m }
 {
   \clist_map_inline:nn {#3}
     {
       \exp_args:Nc \cs_if_exist:NTF { #1 ##1 }
         { \msg_error:nnn { pfpl } { alias-duplicate } { #1 ##1 } }
         { \exp_args:Ncc \let { #1 ##1 } { #2 ##1 } }
     }
 }
\ExplSyntaxOff
